\subsection{Angles and Their Numerical Measure}

An angle exists geometrically before it is assigned a number. We can compare,
copy, add, and bisect angles without degrees or radians. A numerical angle
measure requires an additional choice of unit.

Degrees divide one full turn into $360$ equal parts. Radians use the geometry
of a circle. If an angle at the centre of a circle cuts off an arc of length
$s$ on a circle of radius $r$, its radian measure is
\[
\theta=\frac{s}{r}.
\]
The ratio does not depend on the size of the circle: scaling the circle by a
factor scales both $s$ and $r$ by that factor. Again, similarity makes the
quantity well defined.

\begin{remarkbox}
At this point no vector has been introduced. We have only Euclidean objects,
postulates, constructions, similarity, trigonometric ratios, and numerical
angle measure. Vectors will appear only in the next chapter, after the
Cartesian coordinate system has been constructed.
\end{remarkbox}

The next step is to choose two perpendicular lines, an origin, directions, and
a unit length. Those choices will convert the Euclidean plane into a Cartesian
coordinate plane.
